\documentclass[11pt]{article}

\usepackage[T1]{fontenc}
\usepackage{amsmath,amssymb,amsthm}
\usepackage{newpxtext,newpxmath}
\usepackage{mathtools,bm}
\usepackage[margin=1in,headheight=14pt]{geometry}
\usepackage{microtype}
\usepackage{booktabs,longtable,tabularx,array}
\usepackage{enumitem}
\usepackage{xcolor}
\usepackage{fancyhdr}
\usepackage{hyperref}
\usepackage[nameinlink,capitalize,noabbrev]{cleveref}
\usepackage[most]{tcolorbox}
\usepackage{listings}
\usepackage{url}

\definecolor{deepblue}{RGB}{31,67,112}
\definecolor{deepgreen}{RGB}{32,104,80}
\definecolor{warmgray}{RGB}{90,84,76}
\definecolor{paleblue}{RGB}{239,245,251}
\definecolor{palegreen}{RGB}{239,248,244}
\definecolor{paleyellow}{RGB}{253,249,232}
\definecolor{palered}{RGB}{253,240,239}

\hypersetup{
  colorlinks=true,
  linkcolor=deepblue,
  citecolor=deepgreen,
  urlcolor=deepblue,
  pdftitle={Antiderivatives of Monomially Weighted Fabius-Type Functions},
  pdfauthor={OpenAI, research report prepared for Vladimir Reshetnikov}
}

\pagestyle{fancy}
\fancyhf{}
\fancyhead[L]{\small Antiderivatives of Fabius-Type Functions}
\fancyhead[R]{\small \thepage}
\renewcommand{\headrulewidth}{0.4pt}

\setlist[itemize]{leftmargin=1.6em,itemsep=0.25em,topsep=0.4em}
\setlist[enumerate]{leftmargin=1.8em,itemsep=0.3em,topsep=0.4em}
\setlength{\parindent}{1.1em}
\setlength{\parskip}{0.15em}
\allowdisplaybreaks

\newtheorem{theorem}{Theorem}[section]
\newtheorem{proposition}[theorem]{Proposition}
\newtheorem{lemma}[theorem]{Lemma}
\newtheorem{corollary}[theorem]{Corollary}
\newtheorem{conjecture}[theorem]{Conjecture}
\theoremstyle{definition}
\newtheorem{definition}[theorem]{Definition}
\newtheorem{example}[theorem]{Example}
\theoremstyle{remark}
\newtheorem{remark}[theorem]{Remark}

\newtcolorbox{resultbox}[1][]{
  enhanced,breakable,colback=paleblue,colframe=deepblue,
  boxrule=0.7pt,arc=1.5mm,left=1.5mm,right=1.5mm,top=1mm,bottom=1mm,
  title={#1},fonttitle=\bfseries
}
\newtcolorbox{statusbox}[1][]{
  enhanced,breakable,colback=palegreen,colframe=deepgreen,
  boxrule=0.7pt,arc=1.5mm,left=1.5mm,right=1.5mm,top=1mm,bottom=1mm,
  title={#1},fonttitle=\bfseries
}
\newtcolorbox{warningbox}[1][]{
  enhanced,breakable,colback=paleyellow,colframe=warmgray,
  boxrule=0.7pt,arc=1.5mm,left=1.5mm,right=1.5mm,top=1mm,bottom=1mm,
  title={#1},fonttitle=\bfseries
}

\newcommand{\F}{F}
\newcommand{\Fs}{\mathcal F}
\newcommand{\Up}{\operatorname{up}}
\newcommand{\Q}{Q}
\newcommand{\E}{\mathbb E}
\newcommand{\Pp}{\mathbb P}
\newcommand{\R}{\mathbb R}
\newcommand{\C}{\mathbb C}
\newcommand{\N}{\mathbb N}
\newcommand{\dd}{\,\mathrm d}
\newcommand{\e}{\mathrm e}
\newcommand{\ii}{\mathrm i}
\newcommand{\fall}[2]{#1^{\underline{#2}}}
\newcommand{\rise}[2]{#1^{\overline{#2}}}
\newcommand{\Jn}{\mathcal J}
\newcommand{\RL}{\mathcal I}
\newcommand{\M}{\mathcal M}
\newcommand{\eps}{\varepsilon}
\newcommand{\1}{\mathbf 1}
\DeclareMathOperator{\sinc}{sinc}
\DeclareMathOperator{\Li}{Li}
\DeclareMathOperator{\Rea}{Re}
\DeclareMathOperator{\Ima}{Im}

\title{\bfseries Antiderivatives of Monomially Weighted\\[2mm]
Fabius-Type Functions\\[2mm]
\large Exact Volterra Ladders, Newton--Mellin Series, Rvachev Folds,\\
Fractional Orders, and Inverse-Quantile Integrals}
\author{Research report prepared for Vladimir Reshetnikov}
\date{27 August 2026}

\begin{document}
\maketitle

\begin{abstract}
Let \(F\) be the bounded Fabius distribution function, let \(\Up\) be Rvachev's
compactly supported up-function, and let \(Q=F^{-1}\) on \([0,1]\).  This report
constructs normalized antiderivatives of \(x^pF(x)\) and their higher-order
Volterra analogues.  For \(p\in\N_0\), the answer is a finite closed expression
in dyadically contracted Fabius values.  More generally, for every complex
exponent \(\alpha\), every integer order \(n\geq1\), and \(0<x\leq1\), we prove
a locally uniformly and absolutely convergent Newton--Volterra expansion
\[
 \frac1{(n-1)!}\int_0^x(x-t)^{n-1}t^\alpha F(t)\dd t
 =\sum_{r\geq0}(-1)^r\binom{n+r-1}{r}\fall{\alpha}{r}
 x^{\alpha-r}2^{\binom{n+r}{2}}F\!\left(\frac{x}{2^{n+r}}\right).
\]
Endpoint flatness makes this an entire function of \(\alpha\), including all
negative powers; nonnegative integral exponents are exactly the terminating
cases.  At \(x=1\), the coefficients become the Fabius moments and give a
Newton interpolation formula, a two-parameter family of exact moment
identities, and a logarithmic constant
\[
 \int_0^1\frac{F(x)}x\dd x
 =-\E\log X
 =\sum_{n\geq1}\frac{d_n}{n}
 =\log2+\frac12\sum_{n\geq1}\frac{c_n}{n}.
\]
We derive global piecewise formulas, exponential generating functions,
Bell--Bernoulli cumulant forms, formulas for \(F(1-x)=1-F(x)\), complete
piecewise antiderivatives of \(x^p\Up(x)\), exact higher derivatives in
Thue--Morse form, and Riemann--Liouville extensions of nonintegral order.
For the inverse Fabius function we obtain closed change-of-variable formulas,
including
\[
 \int_0^yQ(u)\dd u=yQ(y)-F\!\left(\frac{Q(y)}2\right),\qquad
 \RL^\nu Q(y)=\frac1{\Gamma(\nu+1)}
 \int_0^{Q(y)}(y-F(t))^\nu\dd t.
\]
Combining the inverse endpoint equivalent already developed in the repository
with regular-variation arguments yields new asymptotics for weighted inverse
primitives and for large quantile moments.  The final sections separate proved
statements from conjectural saddle refinements, record reproducible numerical
checks, and propose Lean formalization and arithmetic directions.
\end{abstract}

\tableofcontents
\newpage

\section{Scope, status, and corpus audit}

\subsection{The source corpus}
The requested source tree was audited recursively at the repository snapshot
identified in its own frontier inventory as commit
\texttt{32d6d36c51d803289e6d6a0dc0c37753766eba47} (27 August 2026).  That
inventory records 67 distinct \texttt{*.tex} paths under
\path{Analysis/FabiusFunction/docs}.  The audit included the live primary
exposition, the mathematical glossary, the Lean walkthrough, the
non-elementarity article, the consolidated non-formalized frontier volume,
the source-paper transcriptions, and the frozen archive and draft families
\cite{ProveItPrimary,ProveItFrontiers,ProveItGlossary,ProveItLean,
ProveItNonElementary,ProveItInventory}.

Archived files contain many historical copies of material later consolidated
into the live exposition or frontier volume.  For nonduplication, the present
work therefore used two passes:
\begin{enumerate}
\item a path-complete inventory and phrase audit over the recursive corpus;
\item theorem-level comparison against the live primary exposition and the
      consolidated frontier volume, with targeted checks of the inverse,
      repeated-integration, Newton/Richardson, moment, asymptotic, and
      non-elementarity source families.
\end{enumerate}
The corpus already contains the basic functional equations, exact dyadic
values, moment recurrences, the sinc product, all-orders endpoint asymptotics,
repeated-averaging limits, and extensive inverse-function work.  It does
\emph{not} develop the normalized antiderivative ladder of this report, the
finite monomial formula, the arbitrary-complex-power Newton--Volterra series,
or the inverse-primitive reductions below.  Literal searches for
``antiderivative'' and ``indefinite integral'' returned no such treatment in
the live primary or consolidated frontier source.  This establishes novelty
\emph{relative to the audited repository}; it is not a claim that no equivalent
identity has ever appeared elsewhere.  In particular, the base functional
integral identity \(\int_0^xF(t)\dd t=F(x/2)\) is an immediate classical
consequence of the defining equation.

\begin{statusbox}[Status convention used in this report]
Every item labeled theorem, proposition, lemma, or corollary is either proved
in place, derived immediately from a proved identity, or explicitly identified
as an imported standard theorem.  In particular, the first formula in
\cref{thm:Q-derivatives} is the classical Lagrange inverse-derivative formula;
the report cites that result and proves only the coefficient conversion used
here.  ``Repository input'' means that an already established result is
imported from the audited corpus.  ``Conjecture'' means that the statement is
supported by the preceding analysis but is not proved.  Apart from derivative
identities already present in the audited corpus and a post-snapshot Lean proof
of the inverse endpoint equivalent, the new antiderivative and inverse-integral
theorems in this report have not yet been translated into Lean.
\end{statusbox}

\subsection{What ``closed form'' means here}
The repository proves strong non-elementarity statements for \(F\), \(\Up\),
and \(Q\) \cite{ProveItNonElementary,ProveItFrontiers}.  Consequently an
antiderivative of \(x^pF(x)\) cannot generally be elementary on any interior
interval: if it were, its derivative divided by \(x^p\) on a subinterval not
containing zero would make \(F\) elementary there.  The meaningful native
closed class is therefore the algebra generated by polynomials, rational
constants, and finitely many dyadically scaled Fabius values.  The finite
formula in \cref{thm:global-integer} is closed in precisely that sense.

\section{Normalization and probabilistic infrastructure}

\subsection{The three Fabius-type functions}
We follow the repository's careful distinction between three objects.  The
bounded distribution function is
\[
 F(x)=0\quad(x\leq0),\qquad F(x)=1\quad(x\geq1),
\]
with
\begin{align}
 F(1-x)&=1-F(x), &&x\in\R,                                      \label{eq:F-reflect}\\
 F'(x)&=2F(2x), &&0\leq x\leq\tfrac12.                           \label{eq:F-left-diff}
\end{align}
The compactly supported Rvachev function is
\begin{equation}
 \Up(x)=F(1-|x|),                                                \label{eq:up-def}
\end{equation}
and the signed global extension is
\begin{equation}
 \Fs(x)=\sum_{m\geq0}(-1)^{s_2(m)}\Up(x-2m-1),
 \qquad \Fs'(x)=2\Fs(2x).                                      \label{eq:signed-def}
\end{equation}
Here \(s_2(m)\) is the binary digit sum.  The bounded and signed functions
coincide on \(( -\infty,1]\), but not globally
\cite{ProveItPrimary,AriasArithmetic}.

Let \((U_j)_{j\geq1}\) be independent uniform random variables on \([0,1]\),
and put
\begin{equation}
 X=\sum_{j\geq1}2^{-j}U_j.                                      \label{eq:X-def}
\end{equation}
Then \(F\) is the distribution function of \(X\), and \(X\overset d=1-X\).
The density of \(Y=2X-1\) is \(\Up\).  We use the moments
\begin{equation}
 d_n=\E X^n,
 \qquad
 c_n=\E(2X-1)^{2n}=\int_{-1}^1x^{2n}\Up(x)\dd x.                \label{eq:dc-def}
\end{equation}
The repository establishes
\begin{align}
 d_0&=1,\qquad
 (n+1)(2^n-1)d_n=\sum_{k=0}^{n-1}\binom{n+1}{k}d_k,             \label{eq:d-rec}\\
 d_n&=n!\,2^{\binom n2}F(2^{-n}),                               \label{eq:d-dyadic}\\
 c_n&=\frac{2d_{2n+1}}{2n+1}.                                  \label{eq:c-d}
\end{align}
The moment generating function
\begin{equation}
 G(z)=\E\e^{zX}=\sum_{n\geq0}d_n\frac{z^n}{n!}
 =\prod_{j\geq1}\frac{\e^{z/2^j}-1}{z/2^j}                    \label{eq:G-product}
\end{equation}
satisfies
\begin{equation}
 G(2z)=\frac{\e^z-1}{z}G(z),\qquad G(-z)=\e^{-z}G(z).           \label{eq:G-fe}
\end{equation}
At imaginary arguments, the centered transform
\(\e^{-z}G(2z)\) becomes the familiar infinite sinc product for \(\Up\)
\cite{AriasCompact,Rvachev1990,ProveItPrimary}.

\subsection{Normalized integer and fractional primitives}
For an integer \(n\geq1\), define the normalized Volterra primitive
\begin{equation}
 (\Jn_nf)(x)=\frac1{(n-1)!}\int_0^x(x-t)^{n-1}f(t)\dd t.         \label{eq:J-def}
\end{equation}
It is the unique solution of \(D^nu=f\) whose derivatives of orders
\(0,\dots,n-1\) vanish at zero.  Every other \(n\)-fold antiderivative differs
from it by a polynomial of degree at most \(n-1\).

For \(\Rea\nu>0\), the Riemann--Liouville version is
\begin{equation}
 (\RL^\nu f)(x)=\frac1{\Gamma(\nu)}
 \int_0^x(x-t)^{\nu-1}f(t)\dd t.                               \label{eq:RL-def}
\end{equation}
The notation \(\Jn_n=\RL^n\) will be used when the order is integral.
For a complex exponent \(\alpha\) and positive \(x\), we take
\(x^\alpha=\exp(\alpha\log x)\) with the real logarithm.

\section{The dyadic antiderivative ladder}

\begin{theorem}[Exact normalized primitive ladder]\label{thm:ladder}
For every integer \(n\geq1\) and every \(0\leq x\leq1\),
\begin{equation}
 \boxed{\quad
 (\Jn_nF)(x)=2^{\binom n2}F\!\left(\frac{x}{2^n}\right).
 \quad}                                                        \label{eq:J-ladder}
\end{equation}
For the signed global extension, the same formula holds for every real \(x\):
\begin{equation}
 D^{-n}_{0}\Fs(x)=2^{\binom n2}\Fs\!\left(\frac{x}{2^n}\right),              \label{eq:signed-primitive}
\end{equation}
where \(D^{-n}_{0}\) denotes the antiderivative with zero \((n-1)\)-jet at
zero.
\end{theorem}

\begin{proof}
Set \(H_n(x)=2^{\binom n2}F(x/2^n)\).  Since
\(0\leq x/2^n\leq1/2\), \eqref{eq:F-left-diff} gives
\[
 H_n'(x)
 =2^{\binom n2-n}F'(x/2^n)
 =2^{\binom n2-n+1}F(x/2^{n-1})
 =H_{n-1}(x).
\]
Iterating yields \(H_n^{(n)}=F\).  Infinite flatness of \(F\) at zero makes
the first \(n\) initial data vanish, so uniqueness of the Volterra primitive
proves \eqref{eq:J-ladder}.  The identical calculation with
\(\Fs'(x)=2\Fs(2x)\) has no range restriction and proves
\eqref{eq:signed-primitive}.
\end{proof}

\begin{remark}
The range restriction in \eqref{eq:J-ladder} is essential for the bounded
function: it is \(\Fs\), not the bounded CDF, that obeys the clean dilation
equation on all of \(\R\).  The global bounded formula will instead be
piecewise polynomial in \cref{thm:global-integer}.
\end{remark}

\begin{corollary}[Moment interpretation of the ladder]\label{cor:J-moment}
For \(x\geq0\),
\begin{equation}
 (\Jn_nF)(x)=\frac1{n!}\E (x-X)_+^n.                            \label{eq:J-stoploss}
\end{equation}
In particular,
\begin{equation}
 (\Jn_nF)(1)=\frac{d_n}{n!}
 =2^{\binom n2}F(2^{-n}).                                      \label{eq:J-one}
\end{equation}
\end{corollary}

\begin{proof}
Write \(F(t)=\E\1_{\{X\leq t\}}\), apply Tonelli's theorem, and integrate
\((x-t)^{n-1}\) from \(t=X\) to \(t=x\).  At \(x=1\), symmetry gives
\(\E(1-X)^n=\E X^n=d_n\).
\end{proof}

\section{The Newton--Volterra expansion for arbitrary powers}

\subsection{A general commutator identity}
The key algebraic mechanism does not depend on the Fabius function.

\begin{theorem}[Newton--Volterra commutator]\label{thm:commutator}
Let \(n\geq1\), \(x>0\), and suppose the displayed integrals converge.  Then
for every complex \(\alpha\),
\begin{equation}
 \Jn_n[t^\alpha f(t)](x)
 =\sum_{r=0}^{\infty}(-1)^r\binom{n+r-1}{r}\fall{\alpha}{r}
 x^{\alpha-r}(\Jn_{n+r}f)(x),                                 \label{eq:general-commutator}
\end{equation}
whenever termwise integration is justified.  If \(\alpha=p\in\N_0\), the
series terminates at \(r=p\) and is an algebraic identity for every locally
integrable \(f\):
\begin{equation}
 \boxed{\quad
 \Jn_n[t^pf(t)](x)
 =\sum_{r=0}^{p}(-1)^r\binom{n+r-1}{r}\frac{p!}{(p-r)!}
 x^{p-r}(\Jn_{n+r}f)(x).
 \quad}                                                        \label{eq:finite-commutator}
\end{equation}
\end{theorem}

\begin{proof}
Inside the Volterra kernel, expand about the upper endpoint:
\[
 t^\alpha=x^\alpha\left(1-\frac{x-t}{x}\right)^\alpha
 =\sum_{r\geq0}(-1)^r\frac{\fall{\alpha}{r}}{r!}
 x^{\alpha-r}(x-t)^r.
\]
After integration, the factor multiplying \(\Jn_{n+r}f\) is
\[
 \frac{\Gamma(n+r)}{\Gamma(n)r!}
 =\binom{n+r-1}{r}.
\]
For \(p\in\N_0\), the binomial expansion is finite, so no convergence
argument is needed.
\end{proof}

The first two operator recurrences encoded in \eqref{eq:finite-commutator} are
\begin{align}
 \Jn_n[t f(t)]&=x\Jn_nf-n\Jn_{n+1}f,                            \label{eq:xf-rec}\\
 \Jn_n[t^pf(t)]&=x\Jn_n[t^{p-1}f(t)]-n\Jn_{n+1}[t^{p-1}f(t)].   \label{eq:triangular-rec}
\end{align}
They provide a convenient triangular implementation without symbolic
integration.

\subsection{The complete Fabius formula}

\begin{theorem}[Arbitrary-complex-power Fabius primitive]\label{thm:newton-volterra}
Let \(n\geq1\), \(0<x\leq1\), and \(\alpha\in\C\).  Then the integral
\(\Jn_n[t^\alpha F(t)](x)\) exists, is entire as a function of \(\alpha\), and
has the absolutely and locally uniformly convergent expansion
\begin{equation}
 \boxed{\begin{aligned}
 \Jn_n[t^\alpha F(t)](x)
 ={}&\sum_{r=0}^{\infty}(-1)^r\binom{n+r-1}{r}\fall{\alpha}{r}
 2^{\binom{n+r}{2}}x^{\alpha-r}
 F\!\left(\frac{x}{2^{n+r}}\right).
 \end{aligned}}                                                \label{eq:master-alpha-n}
\end{equation}
For \(\alpha=p\in\N_0\), this terminates and becomes
\begin{equation}
 \boxed{\begin{aligned}
 \Jn_n[t^pF(t)](x)
 ={}&\sum_{r=0}^{p}(-1)^r\binom{n+r-1}{r}\frac{p!}{(p-r)!}
 2^{\binom{n+r}{2}}x^{p-r}
 F\!\left(\frac{x}{2^{n+r}}\right).
 \end{aligned}}                                                \label{eq:finite-p-n}
\end{equation}
\end{theorem}

\begin{proof}
Substitution of \eqref{eq:J-ladder} into
\eqref{eq:general-commutator} gives the formal expression.  It remains to
justify convergence for all \(\alpha\).

Infinite flatness implies that for every integer \(M\geq0\) there is a finite
constant \(C_M\) such that \(F(t)\leq C_Mt^M\) on \([0,1]\).  Hence, for every
integer \(m\geq1\),
\begin{equation}
 0\leq(\Jn_mF)(x)
 \leq \frac{C_MM!}{(m+M)!}x^{m+M}.                             \label{eq:J-flat-bound}
\end{equation}
Indeed, this is the beta integral obtained by replacing \(F(t)\) with
\(C_Mt^M\).  The absolute value of the \(r\)-th term of
\eqref{eq:general-commutator}, with \(m=n+r\), is therefore at most
\begin{equation}
 C_MM!x^{\Rea\alpha+n+M}
 \binom{n+r-1}{r}\frac{|\fall{\alpha}{r}|}{(n+r+M)!}.           \label{eq:term-bound}
\end{equation}
On a compact set of \(\alpha\)'s, the falling factorial is bounded by
\(C\,r!r^B\) for some fixed \(B\).  Since
\(\binom{n+r-1}{r}=O(r^{n-1})\), choosing \(M>B+1\) makes
\eqref{eq:term-bound} summable uniformly on that compact set.  This proves
absolute local uniform convergence.  The same flatness bound, with logarithmic
powers inserted, permits differentiation under the defining integral and
shows that it is entire in \(\alpha\).  The series and integral agree first in
a half-plane where the ordinary binomial argument is immediate and then
everywhere by the identity theorem; equivalently, one may use dominated
convergence directly with \eqref{eq:term-bound}.
\end{proof}

\begin{resultbox}[Main closed form for the ordinary antiderivative]
For \(0\leq x\leq1\) and \(p\in\N_0\), the normalized ordinary primitive is
\begin{equation}
 \int_0^x t^pF(t)\dd t
 =\sum_{r=0}^{p}(-1)^r\frac{p!}{(p-r)!}
 2^{\binom{r+1}{2}}x^{p-r}
 F\!\left(\frac{x}{2^{r+1}}\right).                            \label{eq:ordinary-p}
\end{equation}
Thus only \(p+1\) dyadically contracted Fabius evaluations are required.
\end{resultbox}

A useful one-step recurrence follows by integrating by parts once and changing
variables:
\begin{equation}
 A_p(x):=\int_0^xt^pF(t)\dd t
 =x^pF(x/2)-p\,2^pA_{p-1}(x/2),\qquad A_0(x)=F(x/2).             \label{eq:A-rec}
\end{equation}
Iteration of \eqref{eq:A-rec} is exactly \eqref{eq:ordinary-p}.

\subsection{Remainders and the negative-power series}
For the ordinary primitive \(n=1\), repeated integration by parts gives an
especially transparent exact remainder.  For every integer \(N\geq0\),
\begin{align}
 \int_0^xt^\alpha F(t)\dd t
 ={}&\sum_{r=0}^{N}(-1)^r\fall{\alpha}{r}
 x^{\alpha-r}(\Jn_{r+1}F)(x)+R_N(\alpha,x),                    \label{eq:partial-alpha}\\
 R_N(\alpha,x)
 ={}&(-1)^{N+1}\fall{\alpha}{N+1}
 \int_0^xt^{\alpha-N-1}(\Jn_{N+1}F)(t)\dd t.                  \label{eq:remainder-alpha}
\end{align}
If \(M> -\Rea\alpha-1\), \eqref{eq:J-flat-bound} yields
\begin{equation}
 |R_N(\alpha,x)|
 \leq
 \frac{|\fall{\alpha}{N+1}|C_MM!}
 {(N+M+1)!(\Rea\alpha+M+1)}
 x^{\Rea\alpha+M+1}.                                         \label{eq:remainder-bound}
\end{equation}
This estimate is useful even when direct evaluation of the extremely small
Fabius terms would underflow.

For a negative integer exponent \(\alpha=-m\), \(m\geq1\), the signs cancel
and every term is positive:
\begin{equation}
 \boxed{\begin{aligned}
 \int_0^xt^{-m}F(t)\dd t
 =\sum_{r=0}^{\infty}
 \frac{(m+r-1)!}{(m-1)!}
 2^{\binom{r+1}{2}}x^{-m-r}
 F\!\left(\frac{x}{2^{r+1}}\right).
 \end{aligned}}                                                \label{eq:negative-series}
\end{equation}
No principal value or regularization is involved: the flat zero of \(F\) makes
\(x^{-m}F(x)\) integrable for every \(m\).  This behavior is sharply different
from \(x^{-m}\Up(x)\), which is singular at the origin because \(\Up(0)=1\).

\begin{corollary}[Endpoint localization]\label{cor:endpoint-localization}
Fix \(n\geq1\) and \(\alpha\in\C\) with \(\Rea\alpha\geq0\).  As
\(x\downarrow0\),
\begin{equation}
 \Jn_n[t^\alpha F(t)](x)
 \sim x^\alpha(\Jn_nF)(x)
 =x^\alpha2^{\binom n2}F(x/2^n).                               \label{eq:localization}
\end{equation}
\end{corollary}

\begin{proof}
The sharp endpoint theory in the audited corpus implies rapid variation:
for every fixed \(0<c<1\), \(F(cx)/F(x)\to0\) as \(x\downarrow0\).  Normalize
the nonnegative measure
\((x-t)^{n-1}F(t)\dd t\) on \([0,x]\).  Its mass below \((1-\delta)x\) is
negligible relative to its total mass, because comparison with fixed dilates
of \(F\) shows concentration in \([(1-\delta)x,x]\).  On that interval
\((t/x)^\alpha\to1\) uniformly as \(\delta\downarrow0\).  Moreover,
\(|(t/x)^\alpha|=(t/x)^{\Rea\alpha}\leq1\) on \((0,x]\), so the complementary
part contributes at most twice its vanishing normalized mass.  The quotient of
the two sides of \eqref{eq:localization} is the expectation of
\((t/x)^\alpha\) under this normalized measure, hence tends to one.
\end{proof}

For \(\Rea\alpha<0\), the same localization is expected and is consistent
with the exact Newton--Volterra series, but the preceding concentration
argument alone does not prove it: \((t/x)^\alpha\) is then unbounded near
zero.  Extending \eqref{eq:localization} to that half-plane requires a separate
uniform-integrability estimate from quantitative flatness (or from the sharp
endpoint expansion), and no such extension is asserted here.

The first correction displayed by \eqref{eq:master-alpha-n} is
\begin{equation}
 \frac{\Jn_n[t^\alpha F(t)](x)}{x^\alpha\Jn_nF(x)}
 =1-n\alpha\frac{\Jn_{n+1}F(x)}{x\Jn_nF(x)}+\cdots.             \label{eq:first-local-correction}
\end{equation}
The quotient on the right is governed by the repository's Lambert phase and
periodic saddle fluctuation; extracting its complete expansion is posed as a
frontier problem in \cref{sec:conjectures}.

\section{Global closed form for natural powers}

For integer powers the bounded extension of \(F\) permits a complete all-real
answer.  Define
\begin{equation}
 H_{n,p}(x)=\frac1{(n-1)!}\int_0^x(x-t)^{n-1}t^pF(t)\dd t,       \label{eq:Hnp-def}
\end{equation}
with the usual oriented integral when \(x<0\).

\begin{theorem}[Global piecewise formula]\label{thm:global-integer}
For \(n\geq1\) and \(p\in\N_0\),
\begin{equation}
 H_{n,p}(x)=
 \begin{cases}
  0, & x\leq0,\\[1mm]
  \displaystyle
  \sum_{r=0}^{p}(-1)^r\binom{n+r-1}{r}\frac{p!}{(p-r)!}
  2^{\binom{n+r}{2}}x^{p-r}
  F\!\left(\frac{x}{2^{n+r}}\right),
   & 0\leq x\leq1,\\[4mm]
  \displaystyle
  \frac{p!}{(p+n)!}x^{p+n}
  -\frac1{(n-1)!}\sum_{j=0}^{n-1}(-1)^j\binom{n-1}{j}
  x^{n-1-j}\frac{d_{p+j+1}}{p+j+1},
   & x\geq1.
 \end{cases}                                                    \label{eq:global-Hnp}
\end{equation}
The pieces have matching derivatives through order \(n-1\) at zero and one.
\end{theorem}

\begin{proof}
The first two cases follow from the support convention and
\cref{thm:newton-volterra}.  For \(x\geq1\), write
\[
 H_{n,p}(x)=\Jn_n[t^p](x)
 -\frac1{(n-1)!}\int_0^1(x-t)^{n-1}t^p(1-F(t))\dd t.
\]
The first term is \(p!x^{p+n}/(p+n)!\).  Expand the finite polynomial kernel
and use
\begin{equation}
 \int_0^1t^q(1-F(t))\dd t=\frac{d_{q+1}}{q+1},                 \label{eq:tail-moment}
\end{equation}
which is the standard tail integral for \(\E X^{q+1}\).  Since the formula
was derived from a single Volterra integral, matching jets are automatic.
\end{proof}

For the ordinary primitive, the last branch simplifies to
\begin{equation}
 A_p(x)=\frac{x^{p+1}-d_{p+1}}{p+1},\qquad x\geq1.              \label{eq:Ap-right}
\end{equation}
This makes clear that the normalized antiderivative is constant-polynomial
outside the transition interval, exactly as repeated integration of a compact
transition requires.

\section{Definite integrals, Newton interpolation, and moment identities}

\subsection{The entire Mellin moment}
Endpoint flatness makes
\begin{equation}
 M(s)=\E X^s=\int_0^1x^s\dd F(x)                               \label{eq:M-def}
\end{equation}
an entire function of \(s\in\C\).  Integration by parts, initially in a
right half-plane and then by analytic continuation, gives
\begin{equation}
 \boxed{\quad
 \mathcal A(\alpha):=\int_0^1x^\alpha F(x)\dd x
 =\frac{1-M(\alpha+1)}{\alpha+1}.
 \quad}                                                        \label{eq:A-M}
\end{equation}
The apparent singularity at \(\alpha=-1\) is removable.

Putting \(x=1\) in \eqref{eq:master-alpha-n} and using
\eqref{eq:J-one} gives the central Newton series.

\begin{theorem}[Moment Newton series]\label{thm:moment-newton}
For every \(\alpha\in\C\),
\begin{equation}
 \boxed{\quad
 \mathcal A(\alpha)
 =\sum_{r=0}^{\infty}(-1)^r\binom{\alpha}{r}
 \frac{d_{r+1}}{r+1}.
 \quad}                                                        \label{eq:moment-newton}
\end{equation}
More generally, for every integer \(n\geq1\),
\begin{equation}
 \boxed{\quad
 \Jn_n[t^\alpha F(t)](1)
 =\frac1{(n-1)!}\sum_{r=0}^{\infty}(-1)^r\binom{\alpha}{r}
 \frac{d_{n+r}}{n+r}.
 \quad}                                                        \label{eq:moment-newton-n}
\end{equation}
Both series converge absolutely and locally uniformly in \(\alpha\).
\end{theorem}

The terminology is literal.  If \(\Delta h(z)=h(z+1)-h(z)\), then
\begin{equation}
 \Delta^r\mathcal A(0)
 =\int_0^1(x-1)^rF(x)\dd x
 =(-1)^r\frac{d_{r+1}}{r+1}.                                  \label{eq:Newton-difference}
\end{equation}
Thus \eqref{eq:moment-newton} is the Newton interpolation series of the entire
function \(\mathcal A\) from its nonnegative integral values.

For \(p\in\N_0\), the series terminates and combines with \eqref{eq:A-M} to
give
\begin{equation}
 \boxed{\quad
 \int_0^1x^pF(x)\dd x=\frac{1-d_{p+1}}{p+1}.
 \quad}                                                        \label{eq:integer-definite}
\end{equation}
Together with \eqref{eq:d-dyadic}, this yields the exact finite dyadic identity
\begin{equation}
 \sum_{r=0}^{p}(-1)^r\frac{p!}{(p-r)!}
 2^{\binom{r+1}{2}}F(2^{-r-1})
 =\frac{1-d_{p+1}}{p+1}.                                      \label{eq:dyadic-finite-id}
\end{equation}
Every value in \eqref{eq:dyadic-finite-id} is rational and may be replaced by
the repository's explicit \(q\)-binomial or integer-sequence formulas.

\subsection{A two-parameter finite-difference identity}
Expanding \((1-t)^{n-1}\) in the defining integral at \(x=1\) produces an
independent form.  Equating it with \eqref{eq:moment-newton-n} gives the
following exact identity.

\begin{corollary}[Two-parameter moment identity]\label{cor:two-param}
For \(n\geq1\) and \(p\in\N_0\),
\begin{equation}
 \boxed{\quad
 \sum_{r=0}^{p}(-1)^r\binom pr\frac{d_{n+r}}{n+r}
 =\sum_{j=0}^{n-1}(-1)^j\binom{n-1}{j}
 \frac{1-d_{p+j+1}}{p+j+1}.
 \quad}                                                        \label{eq:two-param-id}
\end{equation}
More generally, for every \(\alpha\in\C\), the finite left-hand side is
replaced by its convergent Newton series and \(d_{p+j+1}\) by
\(M(\alpha+j+1)\):
\begin{equation}
 \sum_{r\geq0}(-1)^r\binom\alpha r\frac{d_{n+r}}{n+r}
 =\sum_{j=0}^{n-1}(-1)^j\binom{n-1}{j}
 \frac{1-M(\alpha+j+1)}{\alpha+j+1},                           \label{eq:two-param-entire}
\end{equation}
where removable values are understood.
\end{corollary}

\subsection{Negative powers and logarithmic constants}
For \(m\geq1\), \eqref{eq:moment-newton} becomes the positive series
\begin{equation}
 \int_0^1x^{-m}F(x)\dd x
 =\sum_{n\geq1}\frac{(m+n-2)!}{(m-1)!\,n!}d_n.                 \label{eq:negative-moment-series}
\end{equation}
For \(m>1\), the Mellin form is
\begin{equation}
 \int_0^1x^{-m}F(x)\dd x=\frac{M(1-m)-1}{m-1}.                \label{eq:negative-M}
\end{equation}
At the removable case \(m=1\),
\begin{equation}
 \boxed{\quad
 \int_0^1\frac{F(x)}x\dd x
 =-M'(0)=-\E\log X
 =\sum_{n\geq1}\frac{d_n}{n}.
 \quad}                                                        \label{eq:log-constant-d}
\end{equation}

The centered moments give a second positive expansion.  Since
\begin{equation}
 M(s)=2^{-s}\sum_{q\geq0}\binom{s}{2q}c_q,                     \label{eq:M-centered}
\end{equation}
and \(\left.\frac{\dd}{\dd s}\binom{s}{2q}\right|_{s=0}
=-1/(2q)\), one obtains
\begin{equation}
 \boxed{\quad
 \int_0^1\frac{F(x)}x\dd x
 =\log2+\frac12\sum_{q\geq1}\frac{c_q}{q}.
 \quad}                                                        \label{eq:log-constant-c}
\end{equation}
Numerically, the accompanying code gives
\begin{equation}
 \int_0^1\frac{F(x)}x\dd x
 =0.75824000040440865\ldots.                                  \label{eq:log-constant-num}
\end{equation}

\section{Exponential master functions and the sinc product}

The entire family of integer-power primitives is encoded by an incomplete
Laplace transform.  Applying \cref{thm:commutator} to \(\e^{zt}\), or simply
expanding \(\e^{-z(x-t)}\) in the Volterra kernel, gives the following.

\begin{theorem}[Exponential Volterra master]\label{thm:exp-master}
For \(n\geq1\), \(0\leq x\leq1\), and \(z\in\C\),
\begin{equation}
 \boxed{\begin{aligned}
 \Jn_n[\e^{zt}F(t)](x)
 ={}&\e^{zx}\sum_{r=0}^{\infty}(-z)^r\binom{n+r-1}{r}
 2^{\binom{n+r}{2}}F\!\left(\frac{x}{2^{n+r}}\right).
 \end{aligned}}                                                \label{eq:exp-master-n}
\end{equation}
The series is entire in \(z\).  At ordinary order \(n=1\),
\begin{equation}
 \int_0^x\e^{zt}F(t)\dd t
 =\e^{zx}\sum_{m\geq1}(-z)^{m-1}2^{\binom m2}F(x/2^m).         \label{eq:exp-master-one}
\end{equation}
At \(x=1\),
\begin{equation}
 \boxed{\quad
 \int_0^1\e^{zt}F(t)\dd t
 =\frac{\e^z-G(z)}{z}
 =\e^z\frac{1-G(-z)}{z},
 \quad}                                                        \label{eq:exp-definite}
\end{equation}
with the removable value \(1/2\) at \(z=0\).
\end{theorem}

Taking \(p\) derivatives with respect to \(z\) at zero recovers
\eqref{eq:finite-p-n}.  Thus \eqref{eq:exp-master-n} is a compact generating
object for all natural monomial weights and all integral antiderivative orders.

For the up-function, define the incomplete transform
\begin{equation}
 U(z,x)=\int_{-1}^x\e^{zt}\Up(t)\dd t.                          \label{eq:U-zx}
\end{equation}
On the left half, translation gives the particularly simple dyadic series
\begin{equation}
 U(z,x)=\e^{zx}\sum_{m\geq1}(-z)^{m-1}2^{\binom m2}
 F\!\left(\frac{x+1}{2^m}\right),\qquad -1\leq x\leq0.        \label{eq:U-left-exp}
\end{equation}
At the right endpoint,
\begin{equation}
 \int_{-1}^1\e^{zt}\Up(t)\dd t
 =\e^{-z}G(2z)=\frac{G(z)-G(-z)}{z}.                            \label{eq:up-mgf}
\end{equation}
For \(z=2\pi\ii\xi\), this is precisely the infinite sinc product, up to the
chosen Fourier normalization.  Consequently all compact up-moments and their
incomplete analogues arise by differentiating one dyadic Fabius series whose
complete endpoint is the sinc product.

\section{Bell and Bernoulli polynomial organization}

Taking logarithms of \eqref{eq:G-product} and using the classical Bernoulli
expansion gives
\begin{equation}
 \log G(z)=\frac z2+
 \sum_{q\geq1}
 \frac{B_{2q}}{2q(2q)!(2^{2q}-1)}z^{2q}.                       \label{eq:logG-Bernoulli}
\end{equation}
Hence the cumulants of \(X\) are
\begin{equation}
 \kappa_1=\frac12,\qquad
 \kappa_{2q}=\frac{B_{2q}}{2q(2^{2q}-1)},\qquad
 \kappa_{2q+1}=0\quad(q\geq1).                                \label{eq:cumulants}
\end{equation}
Writing \(Y_n\) for the complete exponential Bell polynomial,
\begin{equation}
 d_n=Y_n(\kappa_1,\ldots,\kappa_n).                             \label{eq:d-Bell}
\end{equation}
Therefore every definite natural-power primitive has the explicit
Bell--Bernoulli form
\begin{equation}
 \int_0^1x^pF(x)\dd x
 =\frac{1-Y_{p+1}(\kappa_1,\ldots,\kappa_{p+1})}{p+1}.          \label{eq:A-Bell}
\end{equation}
Substitution in \eqref{eq:global-Hnp} similarly expresses every exterior
polynomial coefficient through Bernoulli numbers and Bell polynomials.  This
provides a direct bridge between the antiderivative problem and the
Bell--Bernoulli infrastructure already emphasized in the repository.

\section{Reflection transforms: \texorpdfstring{$F(1-x)$}{F(1-x)} and
\texorpdfstring{$1-F(x)$}{1-F(x)}}

The global reflection identity \eqref{eq:F-reflect} means that the two
transformations requested in the problem are the same function:
\begin{equation}
 F(1-x)=1-F(x),\qquad x\in\R.                                  \label{eq:transforms-same}
\end{equation}
For natural powers and every integer \(n\geq1\), linearity gives the complete
normalized formula
\begin{equation}
 \boxed{\quad
 \Jn_n[x^p(1-F(x))]
 =\frac{p!}{(p+n)!}x^{p+n}-H_{n,p}(x),
 \quad}                                                        \label{eq:complement-np}
\end{equation}
where \(H_{n,p}\) is the global piecewise expression
\eqref{eq:global-Hnp}.  The same formula therefore applies verbatim to
\(x^pF(1-x)\).

For arbitrary complex \(\alpha\) on \(x>0\), an ordinary antiderivative is
\begin{equation}
 \int x^\alpha(1-F(x))\dd x=
 \begin{cases}
 \displaystyle \frac{x^{\alpha+1}}{\alpha+1}
 -\int_0^x t^\alpha F(t)\dd t+C,&\alpha\neq-1,\\[3mm]
 \displaystyle \log x-\int_0^x\frac{F(t)}t\dd t+C,&\alpha=-1,
 \end{cases}                                                    \label{eq:complement-alpha}
\end{equation}
where the Fabius term is evaluated by \eqref{eq:master-alpha-n}.  Unlike the
Fabius integral, a primitive normalized at zero exists only when
\(\Rea\alpha>-1\), because \(1-F(0)=1\).  Formula
\eqref{eq:complement-alpha} remains a valid local antiderivative on every
positive interval avoiding zero for all \(\alpha\).

At the complete endpoint, for \(\Rea\alpha>-1\),
\begin{equation}
 \int_0^1x^\alpha(1-F(x))\dd x
 =\frac{M(\alpha+1)}{\alpha+1}.                                \label{eq:complement-M}
\end{equation}
Together, \eqref{eq:A-M} and \eqref{eq:complement-M} split the elementary
beta integral \(1/(\alpha+1)\) into the lower-CDF and upper-tail contributions.

\section{Rvachev up-function: complete weighted primitives}

\subsection{First-order formulas and the singularity at zero}
On \(0\leq x\leq1\), \(\Up(x)=1-F(x)\), so for
\(\Rea\alpha>-1\),
\begin{equation}
 \int_0^xt^\alpha\Up(t)\dd t
 =\frac{x^{\alpha+1}}{\alpha+1}
 -\sum_{r\geq0}(-1)^r\fall\alpha r
 2^{\binom{r+1}{2}}x^{\alpha-r}
 F\!\left(\frac{x}{2^{r+1}}\right).                            \label{eq:up-positive-alpha}
\end{equation}
The sum is the convergent Fabius series from \eqref{eq:master-alpha-n}.
For \(\alpha=-1\), replace the elementary term by \(\log x\), obtaining a
local primitive but not a convergent integral from zero.

This exposes an important structural contrast:
\begin{itemize}
\item \(x^\alpha F(x)\) is integrable at zero for every
      \(\alpha\in\C\), because \(F\) is infinitely flat;
\item \(x^\alpha\Up(x)\) has the ordinary monomial threshold
      \(\Rea\alpha>-1\), because \(\Up(0)=1\).
\end{itemize}
For a chosen upper-half-plane branch on the negative axis,
\((-u)^\alpha=\e^{\pi\ii\alpha}u^\alpha\), evenness gives, whenever the
integrals converge,
\begin{equation}
 \int_{-1}^xt^\alpha\Up(t)\dd t
 =\e^{\pi\ii\alpha}
 \left[
 \int_0^1u^\alpha\Up(u)\dd u-
 \int_0^{-x}u^\alpha\Up(u)\dd u
 \right],\qquad -1\leq x\leq0.                                \label{eq:up-negative-alpha}
\end{equation}
For integer powers this becomes real and simplifies to finite Fabius sums.

\subsection{All natural powers and all integer orders}
Define the primitive based at the left support endpoint,
\begin{equation}
 U_{n,p}(x)=\frac1{(n-1)!}\int_{-1}^x(x-t)^{n-1}t^p\Up(t)\dd t,
 \qquad n\geq1,\ p\in\N_0.                                    \label{eq:Unp-def}
\end{equation}
Let \(H_{n,j}\) be as in \eqref{eq:Hnp-def}, and define the full up-moments
\begin{equation}
 m_k^{\Up}=\int_{-1}^1t^k\Up(t)\dd t
 =\begin{cases}0,&k\text{ odd},\\ c_{k/2},&k\text{ even}.
 \end{cases}                                                    \label{eq:up-full-moment}
\end{equation}

\begin{theorem}[Global weighted up-primitive]\label{thm:up-global}
For \(n\geq1\) and \(p\in\N_0\),
\begin{equation}
 U_{n,p}(x)=
 \begin{cases}
 0,&x\leq-1,\\[1mm]
 \displaystyle
 \sum_{j=0}^p\binom pj(-1)^{p-j}H_{n,j}(x+1),
 &-1\leq x\leq0,\\[4mm]
 \displaystyle
 \frac{(-1)^p}{(n-1)!}
 \sum_{j=0}^{n-1}\binom{n-1}{j}
 x^{n-1-j}\frac{d_{p+j+1}}{p+j+1}
 +\frac{p!}{(p+n)!}x^{p+n}-H_{n,p}(x),
 &0\leq x\leq1,\\[5mm]
 \displaystyle
 \frac1{(n-1)!}\sum_{j=0}^{n-1}(-1)^j\binom{n-1}{j}
 x^{n-1-j}m_{p+j}^{\Up},
 &x\geq1.
 \end{cases}                                                    \label{eq:up-global}
\end{equation}
Every branch is finite, and the interior branches reduce further to the
scaled-Fabius formula \eqref{eq:finite-p-n}.
\end{theorem}

\begin{proof}
For \(-1\leq x\leq0\), set \(s=t+1\).  Then
\(\Up(t)=F(s)\), \(t^p=(s-1)^p\), and the first finite sum follows.
For \(0\leq x\leq1\), split the integral at zero.  The negative half is
\[
 \frac{(-1)^p}{(n-1)!}\int_0^1(x+1-s)^{n-1}(1-s)^pF(s)\dd s.
\]
Expand the polynomial in \(x\) and use
\[
 \int_0^1(1-s)^qF(s)\dd s=\frac{d_{q+1}}{q+1},
\]
which follows from \cref{cor:J-moment}.  The positive half equals
\(\Jn_n[t^p]-H_{n,p}\).  For \(x\geq1\), the full support has been crossed;
expand \((x-t)^{n-1}\) and insert the moments
\eqref{eq:up-full-moment}.
\end{proof}

The unweighted case is worth displaying separately.  Put \(U_n=U_{n,0}\).
Then
\begin{equation}
 U_n(x)=
 \begin{cases}
 0,&x\leq-1,\\
 \displaystyle 2^{\binom n2}F\!\left(\frac{x+1}{2^n}\right),
 &-1\leq x\leq0,\\[2mm]
 \displaystyle
 \frac1{(n-1)!}\sum_{j=0}^{n-1}\binom{n-1}{j}
 \frac{d_{j+1}}{j+1}x^{n-1-j}
 +\frac{x^n}{n!}
 -2^{\binom n2}F\!\left(\frac{x}{2^n}\right),
 &0\leq x\leq1,\\[4mm]
 \displaystyle
 \frac1{(n-1)!}\sum_{q=0}^{\lfloor(n-1)/2\rfloor}
 \binom{n-1}{2q}c_qx^{n-1-2q},&x\geq1.
 \end{cases}                                                    \label{eq:up-unweighted-n}
\end{equation}
For \(n=1\), this is the cumulative distribution function of \(Y=2X-1\).
For larger \(n\), its right tail is the moment polynomial of the compact up
density.

\section{Exact higher derivatives and inverse differentiation}

\subsection{Thue--Morse derivative formulas}
Let \(\eps_k=(-1)^{s_2(k)}\).  Iterating Rvachev's functional-differential
equation gives the standard exact derivative formula
\begin{equation}
 \boxed{\quad
 \Up^{(n)}(x)=2^{\binom{n+1}{2}}
 \sum_{k=0}^{2^n-1}\eps_k\,
 \Up(2^nx+2^n-1-2k).
 \quad}                                                        \label{eq:up-derivative}
\end{equation}
Using \(F'(x)=2\Up(2x-1)\), one obtains, for \(m\geq1\),
\begin{equation}
 \boxed{\quad
 F^{(m)}(x)=2^{\binom{m+1}{2}}
 \sum_{k=0}^{2^{m-1}-1}\eps_k\,
 \Up(2^mx-1-2k).
 \quad}                                                        \label{eq:F-derivative-up}
\end{equation}
This version is valid globally for the bounded \(F\).  On the shrinking left
cell \(0<x<2^{-m}\), it collapses to
\begin{equation}
 F^{(m)}(x)=2^{\binom{m+1}{2}}F(2^mx),                          \label{eq:F-left-higher}
\end{equation}
while the signed extension satisfies the same dilation formula on all of
\(\R\):
\begin{equation}
 \Fs^{(m)}(x)=2^{\binom{m+1}{2}}\Fs(2^mx).                     \label{eq:Fs-higher}
\end{equation}

Leibniz's rule therefore gives a closed formula for the requested derivatives.
For \(p\in\N_0\),
\begin{equation}
 D^n[x^p\Fs(x)]
 =\sum_{j=0}^{\min(n,p)}\binom nj\frac{p!}{(p-j)!}x^{p-j}
 2^{\binom{n-j+1}{2}}\Fs(2^{n-j}x).                            \label{eq:derivative-signed-p}
\end{equation}
For arbitrary \(\alpha\in\C\) and \(x>0\), replace
\(p!/(p-j)!\) by \(\fall\alpha j\) and sum through \(j=n\).
For the bounded function, use
\begin{equation}
 D^n[x^\alpha F(x)]
 =\sum_{j=0}^{n}\binom nj\fall\alpha j x^{\alpha-j}F^{(n-j)}(x),
                                                                    \label{eq:derivative-bounded-alpha}
\end{equation}
with \eqref{eq:F-derivative-up}.  An identical Leibniz formula combined with
\eqref{eq:up-derivative} handles \(x^\alpha\Up(x)\).

\subsection{The derivative/antiderivative dictionary}
The formulas may be summarized as follows.
\begin{itemize}
\item Positive derivative order introduces a finite Thue--Morse superposition
      of dilated up-functions, or a single dilation for \(\Fs\).
\item Negative integral order on \([0,1]\) introduces a single contracted
      Fabius value, multiplied by \(2^{\binom n2}\).
\item Multiplication by a natural monomial turns either operation into a finite
      triangular sum; a nonintegral monomial changes the antiderivative sum
      from finite to absolutely convergent.
\end{itemize}
This is the exact operator closure that replaces an elementary integration
table for the Fabius family.

\section{Fractional antiderivative order}

The integer ladder has a natural Riemann--Liouville continuation, although the
single-dilation collapse is special to integral order.

\begin{proposition}[Probabilistic fractional primitive]\label{prop:fractional-F}
For \(\Rea\nu>0\) and \(x\geq0\),
\begin{equation}
 \boxed{\quad
 (\RL^\nu F)(x)=\frac1{\Gamma(\nu+1)}\E (x-X)_+^\nu.
 \quad}                                                        \label{eq:fractional-stoploss}
\end{equation}
At \(x=1\), symmetry gives
\begin{equation}
 (\RL^\nu F)(1)=\frac{M(\nu)}{\Gamma(\nu+1)}.                  \label{eq:fractional-one}
\end{equation}
For \(x>1\),
\begin{equation}
 (\RL^\nu F)(x)=\frac{x^\nu}{\Gamma(\nu+1)}
 \sum_{k\geq0}(-1)^k\binom\nu k d_kx^{-k};                    \label{eq:fractional-exterior}
\end{equation}
for \(x=1\), endpoint flatness supplies the limiting interpretation.
\end{proposition}

\begin{proof}
Since \(\Rea\nu>0\), the absolute value of the kernel is
\((x-t)^{\Rea\nu-1}\), which is integrable on the finite triangular region
arising after insertion of
\(F(t)=\E\1_{\{X\leq t\}}\).  Fubini's theorem therefore gives
\eqref{eq:fractional-stoploss}; for positive real \(\nu\), this reduces to the
Tonelli argument in \cref{cor:J-moment}.  At one, replace \(1-X\) by the
identically distributed \(X\).  For \(x>1\), expand
\((1-X/x)^\nu\) and take expectations.
\end{proof}

\begin{theorem}[Fractional-order Newton--Volterra formula]\label{thm:fractional-commutator}
Let \(\Rea\nu>0\), \(0<x\leq1\), and \(\alpha\in\C\).  Then
\begin{equation}
 \boxed{\quad
 \RL^\nu[t^\alpha F(t)](x)
 =\sum_{r\geq0}(-1)^r\frac{\rise\nu r}{r!}\fall\alpha r
 x^{\alpha-r}(\RL^{\nu+r}F)(x).
 \quad}                                                        \label{eq:fractional-alpha}
\end{equation}
The series converges absolutely and locally uniformly in \((\alpha,\nu)\) on
compact subsets of \(\C\times\{\Rea\nu>0\}\).  For
\(p\in\N_0\), it terminates:
\begin{equation}
 \RL^\nu[t^pF(t)](x)
 =\sum_{r=0}^p(-1)^r\binom{\nu+r-1}{r}\frac{p!}{(p-r)!}
 x^{p-r}(\RL^{\nu+r}F)(x).                                    \label{eq:fractional-p}
\end{equation}
At \(x=1\),
\begin{equation}
 \boxed{\quad
 \RL^\nu[t^\alpha F(t)](1)
 =\frac1{\Gamma(\nu)}
 \sum_{r\geq0}(-1)^r\binom\alpha r
 \frac{M(\nu+r)}{\nu+r}.
 \quad}                                                        \label{eq:fractional-one-alpha}
\end{equation}
\end{theorem}

\begin{proof}
Expand \(t^\alpha\) about \(x\) exactly as in
\cref{thm:commutator}.  Raising the kernel power from \(\nu-1\) to
\(\nu+r-1\) contributes
\(\Gamma(\nu+r)/\Gamma(\nu)=\rise\nu r\).  The flatness estimate used in
\cref{thm:newton-volterra}, with gamma functions replacing factorials, proves
normal convergence.  Formula \eqref{eq:fractional-one-alpha} follows from
\eqref{eq:fractional-one} and
\[
 \frac{\rise\nu r}{r!}\fall\alpha r
 \frac1{\Gamma(\nu+r+1)}
 =\frac1{\Gamma(\nu)}\binom\alpha r\frac1{\nu+r}.
\]
\end{proof}

This fractional identity is an exact analytic continuation of the complete
integer-order table.  It also shows that the moment function \(M(s)\), rather
than a new special function, is the natural endpoint datum for fractional
Fabius integration.

\section{Antiderivatives of the inverse Fabius function}

\subsection{Exact first-order reductions}
Let
\begin{equation}
 Q(y)=F^{-1}(y),\qquad 0\leq y\leq1.                            \label{eq:Q-def}
\end{equation}
The repository establishes strict monotonicity of \(F\), smoothness of \(Q\)
on \((0,1)\), reflection
\begin{equation}
 Q(1-y)=1-Q(y),                                                \label{eq:Q-reflect}
\end{equation}
and infinite endpoint slope
\cite{ProveItPrimary,ProveItFrontiers}.  Substitution \(u=F(t)\) yields an
unexpectedly compact primitive.

\begin{theorem}[Closed primitive of the inverse]\label{thm:inverse-first}
For \(0\leq y\leq1\),
\begin{equation}
 \boxed{\quad
 \int_0^yQ(u)\dd u
 =yQ(y)-F\!\left(\frac{Q(y)}2\right).
 \quad}                                                        \label{eq:inverse-first}
\end{equation}
In particular,
\begin{equation}
 \int_0^1Q(u)\dd u=\frac12.                                   \label{eq:inverse-mean}
\end{equation}
\end{theorem}

\begin{proof}
Put \(x=Q(y)\).  Then
\[
 \int_0^yQ(u)\dd u=\int_0^xtF'(t)\dd t
 =xF(x)-\int_0^xF(t)\dd t.
\]
Apply \(F(x)=y\) and the base ladder
\(\int_0^xF(t)\dd t=F(x/2)\).  At \(y=1\), use
\(Q(1)=1\) and \(F(1/2)=1/2\).
\end{proof}

For a complex weight \(\alpha\), the same substitution and one integration by
parts give the following exact reduction to a power of \(F\).

\begin{proposition}[Weighted inverse primitive]\label{prop:inverse-weighted}
Let \(0<y\leq1\) and \(x=Q(y)\).  If \(\alpha\neq-1\) and
\(\Rea\alpha\geq-1\), then
\begin{equation}
 \boxed{\quad
 \int_0^yu^\alpha Q(u)\dd u
 =\frac{y^{\alpha+1}Q(y)}{\alpha+1}
 -\frac1{\alpha+1}\int_0^{Q(y)}F(t)^{\alpha+1}\dd t.
 \quad}                                                        \label{eq:inverse-weighted}
\end{equation}
At the exceptional exponent,
\begin{equation}
 \boxed{\quad
 \int_0^y\frac{Q(u)}u\dd u
 =Q(y)\log y-\int_0^{Q(y)}\log F(t)\dd t.
 \quad}                                                        \label{eq:inverse-log}
\end{equation}
For real \(\alpha<-1\) and every \(y>0\), the improper integral diverges.
At \(y=0\), the convergent cases above extend with value zero by taking the
right-hand limit; this convention is needed because the displayed factors
\(y^{\alpha+1}\) and \(\log y\) are not themselves defined at zero.
\end{proposition}

\begin{proof}
After \(u=F(t)\), the left side is
\(\int_0^xtF(t)^\alpha F'(t)\dd t\).  Use
\(F^\alpha F'=(F^{\alpha+1})'/(\alpha+1)\).  On the boundary line
\(\Rea\alpha=-1\), the endpoint factor still vanishes because of the
multiplying \(t\).  At \(\alpha=-1\), integrate
\(t(\log F(t))'\).  The repository's endpoint asymptotic implies
\(t\log F(t)\to0\).  Finally, \(Q(u)\) decays more slowly than every positive
power of \(u\), so \(u^\alpha Q(u)\) cannot be integrable when
\(\alpha<-1\).
\end{proof}

At \(y=1\) and \(p\in\N_0\),
\begin{equation}
 \boxed{\quad
 \int_0^1u^pQ(u)\dd u
 =\frac1{p+1}\left(1-\int_0^1F(t)^{p+1}\dd t\right).
 \quad}                                                        \label{eq:inverse-definite-p}
\end{equation}
If \(X_1,\ldots,X_{p+1}\) are independent Fabius variables, then
\begin{equation}
 \int_0^1u^pQ(u)\dd u
 =\frac1{p+1}\E\max(X_1,\ldots,X_{p+1}).                       \label{eq:inverse-order-stat}
\end{equation}
Thus weighted quantile integrals are normalized expected maxima.
At \(p=-1\), formula \eqref{eq:inverse-log} gives the convergent constant
\begin{equation}
 \int_0^1\frac{Q(u)}u\dd u=-\int_0^1\log F(t)\dd t.            \label{eq:inverse-log-constant}
\end{equation}

\subsection{Higher and fractional inverse primitives}
The indicator representation
\begin{equation}
 Q(u)=\int_0^1\1_{\{F(t)\leq u\}}\dd t                         \label{eq:Q-indicator}
\end{equation}
turns all higher primitives into one-dimensional integrals of powers of
\(y-F(t)\).

\begin{theorem}[All orders for the inverse]\label{thm:inverse-all-orders}
For \(\Rea\nu>0\) and \(0\leq y\leq1\),
\begin{equation}
 \boxed{\quad
 (\RL^\nu Q)(y)
 =\frac1{\Gamma(\nu+1)}
 \int_0^{Q(y)}(y-F(t))^\nu\dd t.
 \quad}                                                        \label{eq:inverse-fractional}
\end{equation}
For an integer \(n\geq1\), this is
\begin{equation}
 (\Jn_nQ)(y)=\frac1{n!}\int_0^{Q(y)}(y-F(t))^n\dd t.           \label{eq:inverse-integer-n}
\end{equation}
At the complete endpoint,
\begin{equation}
 (\RL^\nu Q)(1)
 =\frac1{\Gamma(\nu+1)}\int_0^1F(t)^\nu\dd t,                 \label{eq:inverse-fractional-one}
\end{equation}
where reflection was used in the last step.
\end{theorem}

\begin{proof}
Insert \eqref{eq:Q-indicator} into \eqref{eq:RL-def}.  Because
\(\Rea\nu>0\), the absolute value of the resulting kernel is integrable on
the finite region \(0\leq F(t)\leq u\leq y\); hence Fubini's theorem permits
interchanging the integrals.  For a fixed \(t\leq Q(y)\), the variable \(u\)
ranges from \(F(t)\) to \(y\), and
\[
 \frac1{\Gamma(\nu)}\int_{F(t)}^y(y-u)^{\nu-1}\dd u
 =\frac{(y-F(t))^\nu}{\Gamma(\nu+1)}.
\]
At \(y=1\), substitute \(t\mapsto1-t\) and use
\(1-F(t)=F(1-t)\).
\end{proof}

For a natural weight \(p\), the same argument gives
\begin{equation}
 \boxed{\quad
 \RL^\nu[u^pQ(u)](y)
 =\frac1{\Gamma(\nu)}\int_0^{Q(y)}
 \int_{F(t)}^y(y-u)^{\nu-1}u^p\dd u\dd t.
 \quad}                                                        \label{eq:inverse-weighted-fractional}
\end{equation}
For \(y>0\), the inner integral is an incomplete beta function:
\begin{equation}
 y^{p+\nu}
 \left[B(p+1,\nu)-B_{F(t)/y}(p+1,\nu)\right].                 \label{eq:inverse-beta}
\end{equation}
When \(p\) is a nonnegative integer and \(\nu=n\) is a positive integer, it
is a finite polynomial in \(y\) and \(F(t)\).

\subsection{Higher derivatives of the inverse}
Let \(x=Q(y)\in(0,1)\).  The classical Lagrange inverse-derivative formula
\cite{Comtet} gives an exact all-orders identity that can be combined with the
Thue--Morse derivative expression \eqref{eq:F-derivative-up}.  We import that
classical identity here rather than reprove Lagrange inversion.

\begin{theorem}[Lagrange formula for inverse derivatives]\label{thm:Q-derivatives}
For \(n\geq1\),
\begin{equation}
 \boxed{\quad
 Q^{(n)}(y)=
 \left.\frac{\dd^{n-1}}{\dd t^{n-1}}
 \left(\frac{t-x}{F(t)-F(x)}\right)^n\right|_{t=x}.
 \quad}                                                        \label{eq:Q-Lagrange}
\end{equation}
If
\begin{equation}
 a_m(x)=\frac{F^{(m+1)}(x)}{(m+1)!F'(x)},                       \label{eq:am-inverse}
\end{equation}
then equivalently
\begin{equation}
 Q^{(n)}(y)=\frac{(n-1)!}{F'(x)^n}
 [z^{n-1}]\left(1+\sum_{m\geq1}a_m(x)z^m\right)^{-n}.          \label{eq:Q-coefficient}
\end{equation}
\end{theorem}

Indeed, writing \(t=x+z\) gives
\[
 \frac{t-x}{F(t)-F(x)}
 =\frac1{F'(x)}
 \left(1+\sum_{m\geq1}a_m(x)z^m\right)^{-1}.
\]
Raising this identity to the \(n\)-th power and taking the
\((n-1)\)-st derivative at \(z=0\) multiplies the coefficient of
\(z^{n-1}\) by \((n-1)!\), proving the equivalence of
\eqref{eq:Q-Lagrange} and \eqref{eq:Q-coefficient}.

The first cases are
\begin{align}
 Q'&=\frac1{F'},                                               \label{eq:Q1}\\
 Q''&=-\frac{F''}{(F')^3},                                    \label{eq:Q2}\\
 Q'''&=\frac{3(F'')^2-F'F'''}{(F')^5},                         \label{eq:Q3}\\
 Q''''&=\frac{-15(F'')^3+10F'F''F'''-(F')^2F''''}{(F')^7},     \label{eq:Q4}
\end{align}
all evaluated at \(x=Q(y)\).  Inserting \eqref{eq:F-derivative-up} makes each
one a finite Thue--Morse expression in translated, dilated up-values.  On a
left dyadic cell, \eqref{eq:F-left-higher} further collapses every derivative
of \(F\) to a single scaled Fabius value.

\section{New inverse-primitive asymptotics}
\label{sec:inverse-asymptotics}

\subsection{Repository input and slow variation}
The consolidated frontier volume derives from the sharp Lambert expansion the
endpoint equivalent
\begin{equation}
 Q(y)\sim
 \frac{\sqrt{\log(1/y)}}{\e\sqrt{\log2}}
 \exp\!\left[-\sqrt{2\log2\,\log(1/y)}\right]
 \qquad(y\downarrow0).                                        \label{eq:Q-asympt-input}
\end{equation}
At the audited snapshot this statement was marked ``Derived'' rather than
Lean-formalized in the source.  The current Lean development now proves
precisely this equivalence as
\path{Fabius.fabiusInv_isEquivalent_fabiusInverseAsymptoticMain} in
\path{FabiusInverseAsymptotic.lean}.  Thus the results below are
unconditional mathematical consequences of a now-formalized input, although
their integral and order-statistic deductions have not themselves yet been
translated into Lean \cite{ProveItFrontiers}.

Put
\begin{equation}
 L=\log2,\qquad a=\sqrt{2L},\qquad K=\frac1{\e\sqrt L}.          \label{eq:Ka}
\end{equation}
Then \(Q(y)\sim K\sqrt T\e^{-a\sqrt T}\), where
\(T=\log(1/y)\).  For every fixed \(c>0\), replacing \(T\) by
\(T-\log c\) changes this expression by a factor tending to one.  Therefore
\begin{equation}
 \frac{Q(cy)}{Q(y)}\longrightarrow1,\qquad y\downarrow0,       \label{eq:Q-slow}
\end{equation}
so \(Q\) is slowly varying at zero.

\subsection{Weighted primitives near zero}

\begin{theorem}[Karamata law for inverse primitives]\label{thm:Q-Karamata}
Let \(\alpha>-1\) and \(\nu>0\) be real.  Then, as \(y\downarrow0\),
\begin{equation}
 \boxed{\quad
 \RL^\nu[u^\alpha Q(u)](y)
 \sim
 \frac{\Gamma(\alpha+1)}{\Gamma(\alpha+\nu+1)}
 y^{\alpha+\nu}Q(y).
 \quad}                                                        \label{eq:Q-weighted-Karamata}
\end{equation}
In particular,
\begin{align}
 \int_0^yu^\alpha Q(u)\dd u
 &\sim\frac{y^{\alpha+1}Q(y)}{\alpha+1},                       \label{eq:Q-first-Karamata}\\
 (\RL^\nu Q)(y)&\sim\frac{y^\nu Q(y)}{\Gamma(\nu+1)}.         \label{eq:Q-fractional-Karamata}
\end{align}
\end{theorem}

\begin{proof}
Scale \(u=ys\) in the defining integral:
\[
 \RL^\nu[u^\alpha Q(u)](y)
 =\frac{y^{\alpha+\nu}Q(y)}{\Gamma(\nu)}
 \int_0^1(1-s)^{\nu-1}s^\alpha\frac{Q(ys)}{Q(y)}\dd s.
\]
Pointwise convergence to one follows from \eqref{eq:Q-slow}.  Potter bounds
for slowly varying functions provide an integrable dominating beta kernel.
The limiting integral is
\(B(\alpha+1,\nu)/\Gamma(\nu)
=\Gamma(\alpha+1)/\Gamma(\alpha+\nu+1)\).
\end{proof}

The boundary exponent \(\alpha=-1\) is convergent but belongs to a different
scale.

\begin{theorem}[Logarithmic inverse primitive]\label{thm:Q-log-asympt}
As \(y\downarrow0\), with \(T=\log(1/y)\),
\begin{align}
 \int_0^y\frac{Q(u)}u\dd u
 &\sim\frac{2\sqrt T}{a}Q(y)                                  \label{eq:Q-log-relative}\\
 &\sim\frac{\sqrt2}{\e\log2}\,T
 \exp\!\left[-\sqrt{2\log2\,T}\right].                     \label{eq:Q-log-explicit}
\end{align}
\end{theorem}

\begin{proof}
With \(u=\e^{-s}\),
\[
 \int_0^y\frac{Q(u)}u\dd u=\int_T^\infty Q(\e^{-s})\dd s.
\]
The input equivalent turns the integrand into
\(K\sqrt s\e^{-a\sqrt s}\).  Substituting \(v=\sqrt s\) gives
\[
 2K\int_{\sqrt T}^{\infty}v^2\e^{-av}\dd v
 =2K\e^{-a\sqrt T}
 \left(\frac{T}{a}+\frac{2\sqrt T}{a^2}+\frac2{a^3}\right).
\]
The leading term is \((2\sqrt T/a)Q(y)\).  A standard asymptotic sandwich
transfers the equivalent from the model integrand to \(Q(\e^{-s})\).
\end{proof}

\subsection{Large weighted quantile moments}
For \(p\in\N_0\), define
\begin{equation}
 q_p=\int_0^1u^pQ(u)\dd u,
 \qquad
 C_p=\int_0^1(1-u)^pQ(u)\dd u.                                 \label{eq:qCp}
\end{equation}
Reflection \eqref{eq:Q-reflect} gives the exact identity
\begin{equation}
 q_p=\frac1{p+1}-C_p.                                          \label{eq:q-reflection}
\end{equation}

\begin{theorem}[Large quantile-moment correction]\label{thm:large-quantile}
As \(p\to\infty\) through the nonnegative integers,
\begin{equation}
 C_p\sim\frac{Q(1/p)}p,                                       \label{eq:Cp-asympt}
\end{equation}
and therefore
\begin{equation}
 \boxed{\quad
 q_p=\frac1{p+1}
 -\frac{\sqrt{\log p}}{\e\sqrt{\log2}\,p}
 \exp\!\left[-\sqrt{2\log2\,\log p}\right](1+o(1)).
 \quad}                                                        \label{eq:qp-asympt}
\end{equation}
\end{theorem}

\begin{proof}
Set \(u=t/p\).  Then
\[
 \frac{pC_p}{Q(1/p)}
 =\int_0^p\left(1-\frac tp\right)^p
 \frac{Q(t/p)}{Q(1/p)}\dd t.
\]
For fixed \(t>0\), the integrand tends to \(\e^{-t}\) by slow variation.
Potter bounds and the elementary exponential bound on the binomial kernel
supply domination after splitting the integral into a compact interval and a
tail.  The limit is \(\int_0^\infty\e^{-t}\dd t=1\).  Substitute
\eqref{eq:Q-asympt-input} at \(y=1/p\) and use
\eqref{eq:q-reflection}.
\end{proof}

Let \(M_n=\max(X_1,\ldots,X_n)\) and
\(m_n=\min(X_1,\ldots,X_n)\).  From
\eqref{eq:inverse-order-stat} and symmetry,
\begin{equation}
 1-\E M_n=\E m_n=\int_0^1F(t)^n\dd t.                           \label{eq:order-stat-id}
\end{equation}
The theorem immediately gives a compact order-statistic equivalent.

\begin{corollary}[Extreme order statistics and high inverse order]\label{cor:extreme-order}
As \(n\to\infty\) through the positive integers,
\begin{equation}
 \boxed{\quad
 \E m_n=1-\E M_n\sim Q(1/n).
 \quad}                                                        \label{eq:extreme-order-asympt}
\end{equation}
Consequently,
\begin{equation}
 (\Jn_nQ)(1)=\frac1{n!}\int_0^1F(t)^n\dd t
 \sim\frac{Q(1/n)}{n!}.                                       \label{eq:high-inverse-order}
\end{equation}
\end{corollary}

The correction \(Q(1/n)\) is stretched-logarithmic: it is smaller than every
negative power of \(\log n\), but larger than \(n^{-\eps}\) for every fixed
\(\eps>0\).  This scale is invisible to a conventional expansion in powers of
\(1/n\).

\section{Numerical experiments and reproducibility}

The archive accompanying this report contains
\path{numerical_experiments.py}, a fully commented Python script using only
the standard library and \texttt{mpmath}.  It computes two independent moment
systems:
\begin{align}
 (n+1)(2^n-1)d_n&=\sum_{k<n}\binom{n+1}{k}d_k,                 \label{eq:num-d}\\
 (4^n-1)c_n&=\sum_{k<n}\binom{2n}{2k}\frac{c_k}{2(n-k)+1}.     \label{eq:num-c}
\end{align}
The second recurrence follows from
\(H(2z)=(\sinh z/z)H(z)\) for the centered moment generating function and is
independent of the first recurrence.  Exact rational arithmetic verifies
\eqref{eq:integer-definite} for \(0\leq p\leq8\) and
\eqref{eq:two-param-id} for \(1\leq n\leq6\), \(0\leq p\leq8\).

\begin{table}[htbp]
\centering
\small
\caption{Independent evaluations of
\(\mathcal A(\alpha)=\int_0^1x^\alpha F(x)\dd x\).  The first column uses the
Newton series \eqref{eq:moment-newton}; the second uses
\eqref{eq:A-M} and the centered expansion \eqref{eq:M-centered}.  The shown
discrepancies are truncation errors, not statistical errors.}
\label{tab:fractional-checks}
\begin{tabular}{@{}rlll@{}}
\toprule
\(\alpha\)&Newton--Fabius series&Centered-moment formula&absolute difference\\
\midrule
\(0.5\)&0.4209681733402671557991&0.4209681733402671557990&\(9.0\times10^{-23}\)\\
\(-0.5\)&0.6073863107431723970836&0.6073863107431723973782&\(2.9\times10^{-19}\)\\
\(-1\)&0.7582400004044086356798&0.7582400004044086568776&\(2.1\times10^{-17}\)\\
\(-2\)&1.3148861440057837837803&1.3148861440058189655548&\(3.5\times10^{-14}\)\\
\(2.3\)&0.2591812728638398877191737152&0.2591812728638398877191737148&\(3.7\times10^{-28}\)\\
\bottomrule
\end{tabular}
\end{table}

The script also checks the exponential master identity at
\(z=0.7+0.2\ii\).  The moment-series and independent infinite-product values
of \(G(z)\) agree to approximately \(24\) decimal places, and the two sides of
\eqref{eq:exp-definite} agree to the same precision.  Full machine-readable
values are in \path{verification_output.txt} and
\path{fractional_power_checks.csv}.

\begin{warningbox}[Numerical interpretation]
The positive series for negative powers converge more slowly as the negative
exponent grows.  The displayed digits for \(\alpha=-2\) are therefore fewer
than for positive or fractional exponents.  The code reports both truncations
and their discrepancy rather than presenting unsupported extra digits.
\end{warningbox}

\section{Conjectures and frontier directions}
\label{sec:conjectures}

\subsection{A periodic saddle expansion for weighted primitives}
The exact Newton--Volterra series and the repository's all-orders periodic
Lambert expansion should combine into an all-orders local expansion of
\eqref{eq:localization}.

\begin{conjecture}[Weighted primitive saddle hierarchy]\label{conj:weighted-saddle}
For fixed \(n\geq1\) and \(\alpha\in\C\), let \(\lambda_n(x)\) be the exact
Lambert phase associated with \(x/2^n\).  There exist smooth one-periodic
functions \(P_{j,n,\alpha}\) such that, for every \(N\geq1\),
\begin{equation}
 \frac{\Jn_n[t^\alpha F(t)](x)}
 {x^\alpha2^{\binom n2}F(x/2^n)}
 =1+\sum_{j=1}^{N-1}
 \frac{P_{j,n,\alpha}(\lambda_n(x))}{\lambda_n(x)^j}
 +O(\lambda_n(x)^{-N}).                                       \label{eq:conj-weighted-saddle}
\end{equation}
The first coefficient is obtained by inserting the phase expansion of
\(\Jn_{n+1}F/(x\Jn_nF)\) into \eqref{eq:first-local-correction}.
\end{conjecture}

This conjecture is structurally plausible because each successive term of
\eqref{eq:master-alpha-n} shifts the exact phase by approximately one unit,
and the periodic fluctuation is therefore sampled coherently rather than
averaged away.

\subsection{Second-order large-quantile correction}
Ignoring the next periodic inverse correction for a moment, expansion of the
explicit leading model in \cref{thm:large-quantile} gives
\[
 \frac{Q(t/p)}{Q(1/p)}
 =1+\frac{a\log t}{2\sqrt{\log p}}+O\!\left(\frac{1+|\log t|^2}{\log p}\right).
\]
Since \(\int_0^\infty\e^{-t}\log t\dd t=-\gamma\), this suggests the
following conditional refinement.

\begin{conjecture}[Universal first correction, modulo periodic inverse jets]
\label{conj:quantile-second}
If the relative error in \eqref{eq:Q-asympt-input} is uniformly
\(o((\log(1/y))^{-1/2})\) on fixed multiplicative windows, then
\begin{equation}
 \frac{pC_p}{Q(1/p)}
 =1-\frac{\gamma\sqrt{2\log2}}{2\sqrt{\log p}}
 +o((\log p)^{-1/2}).                                         \label{eq:conj-quantile-second}
\end{equation}
In the full problem, the repository's periodic saddle jets may add a bounded
oscillatory term at this same or the next scale; determining the precise
interaction is open.
\end{conjecture}

\subsection{Large negative moments}
The positive moment function \(M(s)\) is entire, so all negative moments
\(N(r)=M(-r)=\E X^{-r}\) exist.  Exact integration by parts gives
\begin{equation}
 N(r)=1+r\int_0^1x^{-r-1}F(x)\dd x.                            \label{eq:negative-moment-exact}
\end{equation}
Inserting the repository's exact-phase expansion and applying Laplace's method
in the Lambert variable leads to the next conjecture.

\begin{conjecture}[Periodic large-negative-moment law]\label{conj:negative-moments}
As real \(r\to+\infty\),
\begin{equation}
 \log M(-r)
 =\frac{\log2}{2}r^2-r\log r
 +\left(1+\frac{\log2}{2}\right)r
 +\frac12\log r+C+\Omega(r)+o(1),                              \label{eq:conj-negative-moment}
\end{equation}
where \(C\) is constant and \(\Omega\) is a bounded smooth one-periodic
function determined by the saddle fluctuation \(\Psi\).  For integral \(r\),
the periodic phase is sampled at a fixed saddle class and may collapse into
the constant term.
\end{conjecture}

The quadratic coefficient and the terms through \(\frac12\log r\) follow from
the stationary phase near \(\lambda=r+1/2\); the exact periodic saddle shift
must still be controlled.

\subsection{Holonomicity, arithmetic, and formalization problems}
The following directions appear especially promising.
\begin{enumerate}
\item \textbf{Lean formalization.}  Formalize \cref{thm:ladder} first, then the
finite commutator \eqref{eq:finite-commutator}, the global formula
\eqref{eq:global-Hnp}, and the inverse identity \eqref{eq:inverse-first}.
These need only finite sums, elementary Volterra calculus, and results already
present in the repository.  The normal-convergence proof of
\eqref{eq:master-alpha-n} is a natural second wave.

\item \textbf{Dyadic arithmetic.}  Substitute the repository's exact
\(q\)-binomial evaluator into \eqref{eq:finite-p-n}.  Determine sharp
2-adic valuations and odd denominator factors of
\(H_{n,p}(a/2^m)\).  The two-parameter identity
\eqref{eq:two-param-id} may create cancellations not visible term by term.

\item \textbf{Newton-series uniqueness.}  Establish optimal growth classes in
which \eqref{eq:moment-newton} is the unique Newton interpolation of
\((1-d_{p+1})/(p+1)\).  The superpolynomial decay of \(d_n\) should yield
stronger convergence than generic Carlson-type hypotheses.

\item \textbf{Non-holonomicity.}  Test the conjecture that
\(\mathcal A(\alpha)\) and \(M(\alpha)\) are not D-finite and that their
integer samples are not P-recursive.  The dyadic functional equation for
\(G\), the infinite Bernoulli cumulant sequence, and the periodic saddle
fluctuation all point in this direction, but none alone is a proof.

\item \textbf{Fractional order and Mellin--Barnes methods.}  Develop
\eqref{eq:fractional-alpha} as a two-complex-variable special function in
\((\alpha,\nu)\).  Its endpoint form \eqref{eq:fractional-one-alpha} may admit
Mellin--Barnes continuation and residue expansions tied to the sinc product.

\item \textbf{Inverse powers of the CDF.}  Study
\(\int_0^xF(t)^s\dd t\), the sole new object in
\eqref{eq:inverse-weighted}.  For positive integral \(s\) it represents order
statistics; near \(s=0\) its one-sided derivatives involve
\(\int(\log F)^k\), whose growth should encode the quadratic-log endpoint.

\item \textbf{Fourier antiderivatives.}  The complete transform
\eqref{eq:up-mgf} is an infinite sinc product, while the exterior branches of
\eqref{eq:up-global} are moment polynomials.  A distributional Fourier theory
that explicitly separates those polynomial tails from
\((2\pi\ii\xi)^{-n}\widehat{\Up}(\xi)\) should produce useful compact
renormalized primitives.

\item \textbf{Order statistics and periodic corrections.}  Push
\cref{thm:large-quantile} to all orders using the closed periodic inverse jets
already present in the frontier volume.  This would transfer the Lambert--W
and periodic structure into a new probabilistic observable:
\(\E\min(X_1,\ldots,X_n)\).
\end{enumerate}

\section{Summary of principal formulas}

\begin{longtable}{@{}p{0.28\textwidth}p{0.66\textwidth}@{}}
\toprule
Object & Exact formula or leading result\\
\midrule
\endfirsthead
\toprule
Object & Exact formula or leading result\\
\midrule
\endhead
\(n\)-fold primitive of \(F\) on \([0,1]\)&
\(\Jn_nF(x)=2^{\binom n2}F(x/2^n)\).\\[1mm]
Natural monomial weight&
Finite sum \eqref{eq:finite-p-n}; global completion \eqref{eq:global-Hnp}.\\[1mm]
Complex monomial weight&
Absolutely convergent Newton--Volterra series \eqref{eq:master-alpha-n}, entire
in \(\alpha\).\\[1mm]
Complete endpoint&
\(\int_0^1x^\alpha F(x)\dd x=(1-M(\alpha+1))/(\alpha+1)\) and
Newton series \eqref{eq:moment-newton}.\\[1mm]
Negative integer weight&
Positive dyadic series \eqref{eq:negative-series}; logarithmic constant
\eqref{eq:log-constant-d}--\eqref{eq:log-constant-c}.\\[1mm]
Reflection/complement&
\(F(1-x)=1-F(x)\), and
\(\Jn_n[x^p(1-F)]=p!x^{p+n}/(p+n)!-H_{n,p}\).\\[1mm]
Rvachev up-function&
Complete four-branch formula \eqref{eq:up-global}; exterior branches are
moment polynomials.\\[1mm]
Higher derivatives&
Thue--Morse sums \eqref{eq:up-derivative} and
\eqref{eq:F-derivative-up}; single dilation for \(\Fs\).\\[1mm]
Fractional order&
Probability formula \eqref{eq:fractional-stoploss} and commutator
\eqref{eq:fractional-alpha}.\\[1mm]
Inverse first primitive&
\(\int_0^yQ(u)\dd u=yQ(y)-F(Q(y)/2)\), as in
\eqref{eq:inverse-first}.\\[1mm]
Inverse arbitrary order&
\(\RL^\nu Q(y)=\Gamma(\nu+1)^{-1}\int_0^{Q(y)}(y-F(t))^\nu\dd t\).\\[1mm]
Inverse endpoint asymptotic&
Weighted Karamata law \eqref{eq:Q-weighted-Karamata}, logarithmic boundary
\eqref{eq:Q-log-explicit}, and large moment correction \eqref{eq:qp-asympt}.\\
\bottomrule
\end{longtable}

\appendix

\section{Corpus comparison map}

\begin{longtable}{@{}p{0.29\textwidth}p{0.65\textwidth}@{}}
\toprule
Source family & Role in the nonduplication audit\\
\midrule
\endfirsthead
\toprule
Source family & Role in the nonduplication audit\\
\midrule
\endhead
\path{Fabius_Function_and_Rvachev_Up/}&
Primary live exposition: bounded/signed conventions, shape, derivatives,
moments, dyadic arithmetic, effective flatness, sharp Lambert asymptotics, and
formalization status.\\[1mm]
\path{FabiusFunction_Mathematical_Glossary/}&
Notation and cross-reference normalization; used to prevent collisions among
\(F\), \(\Up\), the signed extension, moments, and inverse conventions.\\[1mm]
\path{fabius_lean_walkthrough/}&
Map from human statements to Lean declarations; used to identify which imported
facts are already formalized.\\[1mm]
\path{Non_Elementarity_of_the_Fabius_Function/}&
Analytic-locus and non-elementarity results; establishes why the native finite
Fabius formula, rather than an elementary primitive, is the appropriate notion
of closed form.\\[1mm]
Consolidated frontier volume&
Consolidated frontier volume: repeated averaging, inverse function,
small-argument inversion, q-calculus, sinc-product decay, Lambert--W methods,
periodic saddle jets, and other deferred claims.\\[1mm]
Frontier draft packages&
Package-level checks of inverse, exponent-sequence/Newton, repeated-integration,
Fourier-decay, Thue--Morse, and asymptotic reports.  The Newton report concerns
exponent sequences and q-Richardson reconstruction, not Newton series of
monomial antiderivatives.\\[1mm]
\path{papers/}&
Source transcriptions and corrected papers, especially Arias de Reyna's compact
support and arithmetic articles; used for historical formulas and moment
normalizations.\\[1mm]
\path{archive/}&
Frozen early notes, primary-exposition revisions, q-calculus studies, and
standalone reports.  Historical descendants were treated as lineages rather
than independent novelty claims.\\
\bottomrule
\end{longtable}

The decisive negative checks were not based only on titles.  The live primary
source and consolidated frontier source were searched for the literal phrases
``antiderivative'' and ``indefinite integral''; neither contains a treatment of
the present problem.  The repeated-integration material concerns iterative
averaging/convergence to \(\Up\), not the Volterra primitives here.  The
inverse material treats regularity, dyadic inversion, non-elementarity, and
endpoint asymptotics, but not \eqref{eq:inverse-first},
\eqref{eq:inverse-fractional}, or weighted quantile integrals.  The Newton draft
treats exponent sequences, sinc products, Prouhet cancellation, and
q-Richardson reconstruction, not \eqref{eq:moment-newton}.

\section{Exact initial values and sample antiderivatives}

The first moments computed from \eqref{eq:d-rec} are
\begin{align*}
 d_0&=1,&
 d_1&=\frac12,&
 d_2&=\frac5{18},&
 d_3&=\frac16,\\
 d_4&=\frac{143}{1350},&
 d_5&=\frac{19}{270},&
 d_6&=\frac{1153}{23814},&
 d_7&=\frac{583}{17010}.
\end{align*}
The first centered moments are
\begin{align*}
 c_0&=1,&
 c_1&=\frac19,&
 c_2&=\frac{19}{675},&
 c_3&=\frac{583}{59535},&
 c_4&=\frac{132809}{32531625}.
\end{align*}

For \(0\leq x\leq1\), the first ordinary primitives are
\begin{align}
 \int_0^xF(t)\dd t
 &=F(x/2),                                                     \label{eq:sample-p0}\\
 \int_0^xtF(t)\dd t
 &=xF(x/2)-2F(x/4),                                           \label{eq:sample-p1}\\
 \int_0^xt^2F(t)\dd t
 &=x^2F(x/2)-4xF(x/4)+16F(x/8),                               \label{eq:sample-p2}\\
 \int_0^xt^3F(t)\dd t
 &=x^3F(x/2)-6x^2F(x/4)+48xF(x/8)-384F(x/16).                 \label{eq:sample-p3}
\end{align}
At \(x=1\), these equal respectively
\begin{equation}
 \frac12,\qquad \frac{13}{36},\qquad \frac5{18},\qquad
 \frac{1207}{5400}.                                           \label{eq:sample-definite}
\end{equation}
For a first negative-power example,
\begin{equation}
 \int_0^x\frac{F(t)}t\dd t
 =\sum_{r\geq0}r!\,2^{\binom{r+1}{2}}x^{-r-1}
 F\!\left(\frac{x}{2^{r+1}}\right).                            \label{eq:sample-negative-one}
\end{equation}
For a fractional example,
\begin{equation}
 \int_0^xt^{1/2}F(t)\dd t
 =\sum_{r\geq0}(-1)^r\fall{1/2}{r}
 2^{\binom{r+1}{2}}x^{1/2-r}
 F\!\left(\frac{x}{2^{r+1}}\right).                            \label{eq:sample-half}
\end{equation}

\section{Derivation of the centered-moment recurrence}

Let
\begin{equation}
 H(z)=\E\e^{z(2X-1)}=\sum_{n\geq0}c_n\frac{z^{2n}}{(2n)!}.      \label{eq:H-centered}
\end{equation}
The distributional fixed point
\begin{equation}
 2X-1\ \overset d=\ \frac{V+(2X'-1)}2,                         \label{eq:Y-fixed}
\end{equation}
where \(V\) is uniform on \([-1,1]\) and \(X'\) is an independent copy of
\(X\), gives
\begin{equation}
 H(2z)=\frac{\sinh z}{z}H(z).                                  \label{eq:H-functional}
\end{equation}
Equating coefficients of \(z^{2n}\) yields
\begin{equation}
 (4^n-1)c_n
 =\sum_{k=0}^{n-1}\binom{2n}{2k}\frac{c_k}{2(n-k)+1},          \label{eq:c-recurrence}
\end{equation}
which is the independent recurrence used in the numerical checks.

\section{Implementation notes}

The supplied Python program has four layers.
\begin{enumerate}
\item Exact \texttt{Fraction} arithmetic generates the initial \(d_n\) and
      \(c_n\), checks the terminating natural-power identity, and verifies the
      two-parameter identity with no rounding.
\item High-precision \texttt{mpmath} recurrences generate longer moment arrays.
\item The Newton series \eqref{eq:moment-newton} and centered formula
      \eqref{eq:M-centered} are evaluated independently for positive,
      fractional, negative, and complex exponents.
\item The moment series for \(G\) is compared with the independent product
      \eqref{eq:G-product}, after which \eqref{eq:exp-definite} is checked.
\end{enumerate}
No Monte Carlo sampling is used.  The default parameters in the packaged
output use 85 decimal digits, \(d_0,\ldots,d_{1500}\), and
\(c_0,\ldots,c_{700}\).

\begin{thebibliography}{99}

\bibitem{ProveItPrimary}
V.~Reshetnikov,
\emph{The Fabius Function and Rvachev's Up-Function},
ProveIt repository, live LaTeX exposition, snapshot 27 August 2026.
\url{https://github.com/VladimirReshetnikov/ProveIt/blob/main/Analysis/FabiusFunction/docs/Fabius_Function_and_Rvachev_Up/Fabius_Function_and_Rvachev_Up.tex}.

\bibitem{ProveItFrontiers}
V.~Reshetnikov,
\emph{Non-formalized Research Frontiers for the Fabius--Rvachev System},
consolidated LaTeX volume, ProveIt repository, snapshot 27 August 2026.
\url{https://github.com/VladimirReshetnikov/ProveIt/blob/main/Analysis/FabiusFunction/docs/non-formalized-research-frontiers/non-formalized-research-frontiers.tex}.

\bibitem{ProveItGlossary}
V.~Reshetnikov,
\emph{Fabius Function Mathematical Glossary}, ProveIt repository.
\url{https://github.com/VladimirReshetnikov/ProveIt/tree/main/Analysis/FabiusFunction/docs/FabiusFunction_Mathematical_Glossary}.

\bibitem{ProveItLean}
V.~Reshetnikov,
\emph{Fabius Lean Walkthrough}, ProveIt repository.
\url{https://github.com/VladimirReshetnikov/ProveIt/tree/main/Analysis/FabiusFunction/docs/fabius_lean_walkthrough}.

\bibitem{ProveItNonElementary}
V.~Reshetnikov,
\emph{Non-Elementarity of the Fabius Function}, ProveIt repository.
\url{https://github.com/VladimirReshetnikov/ProveIt/tree/main/Analysis/FabiusFunction/docs/Non_Elementarity_of_the_Fabius_Function}.

\bibitem{ProveItInventory}
V.~Reshetnikov,
\emph{Exponent-Sequence and Newton-Basis Frontiers for the Fabius--Rvachev System},
archive README recording the 67-path recursive LaTeX audit at commit
\texttt{32d6d36c51d803289e6d6a0dc0c37753766eba47}, 27 August 2026.
\url{https://github.com/VladimirReshetnikov/ProveIt/tree/main/Analysis/FabiusFunction/docs/non-formalized-research-frontiers/drafts/Fabius_Newton_Rvachev_Frontier_Report}.

\bibitem{JessenWintner}
B.~Jessen and A.~Wintner,
``Distribution functions and the Riemann zeta function,''
\emph{Transactions of the American Mathematical Society} \textbf{38} (1935),
48--88.  DOI: 10.2307/1989728.

\bibitem{Fabius1966}
J.~Fabius,
``A probabilistic example of a nowhere analytic \(C^\infty\)-function,''
\emph{Zeitschrift f\"ur Wahrscheinlichkeitstheorie und Verwandte Gebiete}
\textbf{5} (1966), 173--174.  DOI: 10.1007/BF00536652.

\bibitem{Rvachev1990}
V.~A.~Rvachev,
``Compactly supported solutions of functional-differential equations and their
applications,'' \emph{Russian Mathematical Surveys} \textbf{45}:1 (1990),
87--120.  DOI: 10.1070/RM1990v045n01ABEH002324.

\bibitem{AriasCompact}
J.~Arias de Reyna,
\emph{An Infinitely Differentiable Function with Compact Support: Definition
and Properties}, arXiv:1702.05442 (2017); English translation of
\emph{Rev. Real Acad. Ciencias Madrid} \textbf{76} (1982), 21--38.

\bibitem{AriasArithmetic}
J.~Arias de Reyna,
``Arithmetic of the Fabius function,'' \emph{Integers} \textbf{18} (2018),
Paper A51; arXiv:1702.06487.

\bibitem{Haugland}
J.~K.~Haugland,
\emph{Evaluating the Fabius Function}, arXiv:1609.07999 (2016).

\bibitem{AlloucheShallit}
J.-P.~Allouche and J.~Shallit,
\emph{Automatic Sequences: Theory, Applications, Generalizations},
Cambridge University Press, 2003.

\bibitem{CorlessLambert}
R.~M.~Corless, G.~H.~Gonnet, D.~E.~G.~Hare, D.~J.~Jeffrey, and D.~E.~Knuth,
``On the Lambert W function,'' \emph{Advances in Computational Mathematics}
\textbf{5} (1996), 329--359.

\bibitem{BinghamRegularVariation}
N.~H.~Bingham, C.~M.~Goldie, and J.~L.~Teugels,
\emph{Regular Variation}, Cambridge University Press, 1987.

\bibitem{SamkoFractional}
S.~G.~Samko, A.~A.~Kilbas, and O.~I.~Marichev,
\emph{Fractional Integrals and Derivatives: Theory and Applications},
Gordon and Breach, 1993.

\bibitem{Comtet}
L.~Comtet,
\emph{Advanced Combinatorics}, Reidel, 1974.

\end{thebibliography}

\end{document}
